\section{Discusión matemática e interpretación}\label{sec:discusion}
La dependencia de $\RRMSE$ respecto a estación removida, distribución de puntos y $\lambda$ muestra la estructura esencial del problema inverso: numerosas combinaciones de coeficientes pueden reproducir observaciones próximas a los sensores y divergir entre ellos. La reconstrucción IDW utiliza geometría local y carece de un vector paramétrico explícito comparable a la SVD; la RBF puede producir un campo más suave, pero su estabilidad depende del ancho de base y penalización. La comparación no establece dominancia universal de un método.

\subsection{Incertidumbre y análisis de sensibilidad}\label{sec:incertidumbre}
Los ensayos de variación de ruido y estaciones están habilitados en interfaces. Sin embargo, no se ha calibrado una distribución de errores instrumental, ni se dispone de residuos reales para estimar una covarianza espacial. Una banda de percentiles de bootstrap por estaciones, cuando se implemente, describiría variabilidad bajo un mecanismo de remuestreo especificado, no una garantía de cobertura. Conviene separar: error de medición en $y_i$; incertidumbre paramétrica en $D,u,\lambda$; error de aproximación entre familia de interpolantes y campo verdadero; y variabilidad espacial/temporal real. Una distribución agrupada induce áreas alejadas sin observaciones y puede hacer engañosa una barra de error basada únicamente en ruido iid.

\subsection{Asociación, predicción y causalidad}
La concentración estimada y los vientos, en el futuro, podrían presentar asociación en zonas del mapa. Tal asociación no identifica por sí sola origen de emisiones porque la concentración integra procesos históricos, química, deposición, mezcla vertical, meteorología y sesgos de observación. Un modelo de transporte basado en \eqref{eq:edp} exige emisiones, parámetros, meteorología, frontera y calibración independiente antes de poder afirmarse como pronóstico urbano. La simulación periódica de este trabajo es únicamente un ensayo de verificación numérica.

\section{Limitaciones del prototipo}
\begin{enumerate}
\item \textbf{Observaciones auténticas ausentes:} los catálogos y endpoints están identificados, pero no se obtuvieron bytes de las series ambientales. No existen porcentajes de datos faltantes calculados ni medidas de cobertura verificadas sobre archivos descargados.
\item \textbf{Dependencias sin ejecución:} MATLAB y R no están disponibles en el entorno actual. Sus módulos constituyen código fuente entregado pendiente de verificación en esas plataformas.
\item \textbf{Dominio sintético:} los puntos no corresponden a ubicaciones de estaciones de Monterrey; la malla no representa límites municipales ni relieve. Las coordenadas locales no son UTM georreferenciadas.
\item \textbf{Hipótesis simplificadas:} los campos sintéticos son suaves, el ruido iid es gaussianamente impuesto y la difusión es uniforme. No hay correlación temporal de medida ni modelos de estaciones con datos faltantes extensos.
\item \textbf{Métodos incompletos:} aún faltan kriging calibrado, splines, gradientes a partir de viento observado, PCA aplicado a datos reales, bootstrap espacial, validación temporal bloqueada con registros auténticos y exportación conjunta de las aplicaciones.
\item \textbf{Interfaces no equivalentes:} el navegador usa \emph{kernel ridge}; MATLAB y Python usan penalización sobre coeficientes por mínimos cuadrados. Por ello sus valores individuales no se deben comparar como si fueran la misma variante RBF.
\item \textbf{Selección de hiperparámetros:} no se ejecutó validación cruzada anidada para seleccionar el mejor $\lambda$ y estimar de manera externa su desempeño.
\end{enumerate}

\section{Conclusiones}\label{sec:conclusion}
Se dispone de un núcleo reproducible para reconstrucción de campos sintéticos, diagnóstico singular y EDP conservativa con pruebas ejecutadas en Python. La regularización es un componente útil cuando el diseño numérico presenta grandes números de condición y debe seleccionarse mediante un protocolo que prevenga fugas de datos. El esquema de volúmenes finitos conserva masa bajo las condiciones periódicas especificadas y reduce su desviación respecto a la solución analítica al refinar la malla. Estas conclusiones son \emph{matemático-computacionales} y se limitan al ensayo declarado. Ningún resultado de este trabajo permite inferir tendencias, peligros, concentraciones reales o causalidad ambiental en Monterrey.

\section{Trabajo futuro}
La secuencia prioritaria consiste en: (i) adquirir y archivar registros SIMA/INECC con metadatos de licencias y unidades; (ii) verificar calidad, coordenadas y horarios contra documentaciones primarias; (iii) desarrollar un módulo de interpolación georreferenciada con máscara urbana/proyección métrica; (iv) realizar validación bloqueada espacial y temporal, también para comparar con líneas base; (v) incorporar un modelo espacial probabilístico con evaluación de calibración; y (vi) utilizar meteorología y fuentes inventariadas para evaluar un modelo físico simplificado, calibrarlo y validar predicciones separadamente. Solo después de estas fases correspondería comunicar conclusiones empíricas sobre el área metropolitana.
